\documentclass[a4paper]{article}
\usepackage[affil-it]{authblk}
\usepackage{amsmath, amssymb}
\usepackage[backend=bibtex,style=numeric]{biblatex}

\usepackage{geometry}
\geometry{margin=1.5cm, vmargin={0pt,1cm}}
\setlength{\topmargin}{-1cm}
\setlength{\paperheight}{29.7cm}
\setlength{\textheight}{25.3cm}

\addbibresource{citation.bib}

\begin{document}
% =================================================
\title{Numerical Analysis homework 2 }

\author{Zhang Zihao 3230101685
  \thanks{Electronic address: \texttt{3230101685@zju.edu.cn}}}
\affil{(Information and Computing Science 1), Zhejiang University }


\date{Due time: \today}

\maketitle

\begin{abstract}
    My homework for numerical analysis.     
\end{abstract}

% ============================================
\begin{document}

\section*{Problem I: Linear Interpolation Error Analysis}
\subsection*{Problem Statement}
For \( f \in \mathcal{C}^2[x_0, x_1] \) and \( x \in (x_0, x_1) \), linear interpolation of \( f \) at \( x_0 \) and \( x_1 \) yields:

\[
f(x) - p_1(f; x) = \frac{f''(\xi(x))}{2}(x - x_0)(x - x_1).
\]

Consider the case \( f(x) = \frac{1}{x} \), \( x_0 = 1 \), \( x_1 = 2 \).

\begin{enumerate}
\item Determine \( \xi(x) \) explicitly.
\item Extend the domain of \( \xi \) continuously from \((x_0, x_1)\) to \([x_0, x_1]\). Find \(\max \xi(x)\), \(\min \xi(x)\), and \(\max f''(\xi(x))\).
\end{enumerate}

\subsection*{Solution}

\paragraph*{Given function and derivatives:}
\[
f(x) = \frac{1}{x}, \quad f'(x) = -\frac{1}{x^2}, \quad f''(x) = \frac{2}{x^3}
\]

\paragraph*{Linear interpolation polynomial:}
The linear interpolant through points \((1, f(1)) = (1, 1)\) and \((2, f(2)) = (2, \frac{1}{2})\) is:
\[
p_1(f; x) = 1 + \frac{\frac{1}{2} - 1}{2 - 1}(x - 1) = 1 - \frac{1}{2}(x - 1) = \frac{3 - x}{2}
\]

\paragraph*{Actual interpolation error:}
\[
f(x) - p_1(f; x) = \frac{1}{x} - \frac{3 - x}{2} = \frac{2 - (3 - x)x}{2x} = \frac{2 - 3x + x^2}{2x} = \frac{(x - 1)(x - 2)}{2x}
\]

\paragraph*{(a) Determining \( \xi(x) \) explicitly:}
From the error formula:
\[
f(x) - p_1(f; x) = \frac{f''(\xi(x))}{2}(x - 1)(x - 2)
\]
Substitute \( f''(\xi(x)) = \frac{2}{\xi(x)^3} \) and the actual error:
\[
\frac{(x - 1)(x - 2)}{2x} = \frac{1}{\xi(x)^3}(x - 1)(x - 2)
\]
For \( x \in (1, 2) \), we can cancel \( (x - 1)(x - 2) \neq 0 \):
\[
\frac{1}{2x} = \frac{1}{\xi(x)^3} \quad \Rightarrow \quad \xi(x)^3 = 2x \quad \Rightarrow \quad \xi(x) = (2x)^{1/3}
\]

\paragraph*{(b) Extension and extremal values:}
Extend \( \xi(x) = (2x)^{1/3} \) continuously to \([1, 2]\):

\begin{itemize}
\item Minimum: \( \min_{x \in [1, 2]} \xi(x) = \xi(1) = (2 \cdot 1)^{1/3} = \sqrt[3]{2} \)
\item Maximum: \( \max_{x \in [1, 2]} \xi(x) = \xi(2) = (2 \cdot 2)^{1/3} = \sqrt[3]{4} \)
\end{itemize}

For \( f''(\xi(x)) = \frac{2}{\xi(x)^3} \):
\[
f''(\xi(x)) = \frac{2}{((2x)^{1/3})^3} = \frac{2}{2x} = \frac{1}{x}
\]
Thus:
\[
\max_{x \in [1, 2]} f''(\xi(x)) = \max_{x \in [1, 2]} \frac{1}{x} = \frac{1}{1} = 1
\]

\subsection*{Final Answers:}
\begin{enumerate}
\item[(a)] \( \xi(x) = (2x)^{1/3} \)
\item[(b)] \( \min \xi(x) = \sqrt[3]{2} \), \( \max \xi(x) = \sqrt[3]{4} \), \( \max f''(\xi(x)) = 1 \)
\end{enumerate}


\section*{Problem II: Non-negative Polynomial Interpolation}

\subsection*{Problem Statement}
Let $\mathbb{P}_m^+$ be the set of all polynomials of degree $\leq m$ that are non-negative on the real line,

\[
\mathbb{P}_m^+ = \{p : p \in \mathbb{P}_m, \, \forall x \in \mathbb{R}, \, p(x) \geq 0\}.
\]

Find $p \in \mathbb{P}_{2n}^+$ such that $p(x_i) = f_i$ for $i = 0, 1, \ldots, n$ where $f_i \geq 0$ and $x_i$ are distinct points on $\mathbb{R}$.

\subsection*{Solution}

Let $q(x)$ be the Lagrange interpolation polynomial of degree $\leq n$ that satisfies:
\[
q(x_i) = \sqrt{f_i}, \quad i = 0, 1, \ldots, n
\]

Explicitly, $q(x)$ can be written as:
\[
q(x) = \sum_{i=0}^n \sqrt{f_i} L_i(x)
\]
where $L_i(x)$ are the Lagrange basis polynomials:
\[
L_i(x) = \prod_{\substack{j=0 \\ j \neq i}}^n \frac{x - x_j}{x_i - x_j}
\]

Now define:
\[
p(x) = [q(x)]^2
\]

\subsection*{Verification of Properties}

\begin{enumerate}
\item \textbf{Degree constraint:} 
Since $q(x) \in \mathbb{P}_n$, we have $p(x) = [q(x)]^2 \in \mathbb{P}_{2n}$.

\item \textbf{Non-negativity:}
For all $x \in \mathbb{R}$, $p(x) = [q(x)]^2 \geq 0$, so $p \in \mathbb{P}_{2n}^+$.

\item \textbf{Interpolation condition:}
For each $x_i$:
\[
p(x_i) = [q(x_i)]^2 = \left[\sqrt{f_i}\right]^2 = f_i
\]
\end{enumerate}

\section*{Problem III: Divided Differences of the Exponential Function}

\subsection*{Problem Statement}
Consider \( f(x) = e^x \).

\begin{enumerate}
\item Prove by induction that  
\[ \forall t \in \mathbb{R}, \quad f[t, t+1, \ldots, t+n] = \frac{(e-1)^n}{n!} e^t. \]

\item From Corollary 2.22 we know  
\[ \exists \xi \in (0, n) \text{ s.t. } f[0, 1, \ldots, n] = \frac{1}{n!} f^{(n)}(\xi). \]

Determine \(\xi\) from the above two equations. Is \(\xi\) located to the left or to the right of the midpoint \(n/2\)?
\end{enumerate}

\subsection*{Solution}

\subsubsection*{Part (a): Proof by Induction}

\textbf{Base case: } \( n = 0 \)

For \( n = 0 \), the divided difference is just the function value:
\[ f[t] = e^t = \frac{(e-1)^0}{0!} e^t = e^t \]
So the formula holds for \( n = 0 \).

\textbf{Inductive step: }

Assume the formula holds for \( n = k \):
\[ f[t, t+1, \ldots, t+k] = \frac{(e-1)^k}{k!} e^t \]

For \( n = k+1 \), using the recurrence relation for divided differences:
\[ f[t, t+1, \ldots, t+k+1] = \frac{f[t+1, \ldots, t+k+1] - f[t, \ldots, t+k]}{(t+k+1) - t} \]

By the inductive hypothesis:
\[
\begin{aligned}
f[t+1, \ldots, t+k+1] &= \frac{(e-1)^k}{k!} e^{t+1}
\end{aligned}
\]

Therefore:
\[
\begin{aligned}
f[t, t+1, \ldots, t+k+1] &= \frac{1}{k+1} \left( \frac{(e-1)^k}{k!} e^{t+1} - \frac{(e-1)^k}{k!} e^t \right) \\
&= \frac{(e-1)^k}{k!(k+1)} e^t (e - 1) \\
&= \frac{(e-1)^{k+1}}{(k+1)!} e^t
\end{aligned}
\]

This completes the inductive step.

\subsubsection*{Part (b): Determining \(\xi\)}

From part (a) with \( t = 0 \):
\[ f[0, 1, \ldots, n] = \frac{(e-1)^n}{n!} \]

From Corollary 2.22:
\[ f[0, 1, \ldots, n] = \frac{1}{n!} f^{(n)}(\xi) = \frac{1}{n!} e^{\xi} \]

Equating both expressions:
\[ \frac{(e-1)^n}{n!} = \frac{1}{n!} e^{\xi} \quad \Rightarrow \quad e^{\xi} = (e-1)^n \]

Taking natural logarithms:
\[ \xi = n \ln(e-1) \]

Since \( e-1 \approx 1.71828 \) and \( \ln(e-1) \approx \ln(1.71828) \approx 0.54132 \), we have:
\[ \xi \approx 0.54132n \]

\subsubsection*{Comparison with midpoint}

The midpoint is \( \frac{n}{2} = 0.5n \). Since:
\[ \xi \approx 0.54132n > 0.5n \]

Therefore, \(\xi\) is located to the \textbf{right} of the midpoint \(n/2\).


\section*{Problem IV: Newton Interpolation and Extremum Location}

\subsection*{Problem Statement}
Given: \( f(0) = 5 \), \( f(1) = 3 \), \( f(3) = 5 \), \( f(4) = 12 \).

\begin{enumerate}
\item Use the Newton formula to obtain \( p_3(f; x) \);
\item The data suggest that \( f \) has a minimum in \( x \in (1, 3) \). Find an approximate value for the location \( x_{\min} \) of the minimum.
\end{enumerate}

\subsection*{Solution}

\subsubsection*{Part 1: Newton Interpolation Polynomial}

Construct the divided difference table:

\[
\renewcommand{\arraystretch}{1.5}
\begin{array}{cccccc}
x_i & f[x_i] & \text{1st order} & \text{2nd order} & \text{3rd order} \\
\hline
0 & 5 \\
 & & f[0,1] = -2 \\
1 & 3 & & f[0,1,3] = 1 \\
 & & f[1,3] = 1 & & f[0,1,3,4] = 0.25 \\
3 & 5 & & f[1,3,4] = 2 \\
 & & f[3,4] = 7 \\
4 & 12 \\
\end{array}
\]

The Newton interpolation polynomial is:
\[
\begin{aligned}
p_3(f; x) &= f[0] + f[0,1](x-0) + f[0,1,3](x-0)(x-1) + f[0,1,3,4](x-0)(x-1)(x-3) \\
&= 0.25x^3 - 2.25x + 5
\end{aligned}
\]

\subsubsection*{Part 2: Approximate Minimum Location}

Since \( p_3(f; x) = 0.25x^3 - 2.25x + 5 \), we find its derivative:
\[
\begin{aligned}
p_3'(x) &= 0.75x^2 - 2.25
\end{aligned}
\]

Set derivative to zero:
\[
\begin{aligned}
0.75x^2 - 2.25 &= 0 \\
x^2 &= 3 \\
x &= \pm \sqrt{3}
\end{aligned}
\]

Since we're looking for a minimum in \( (1, 3) \), we take \( x = \sqrt{3} \approx 1.732 \).

Check second derivative:
\[
\begin{aligned}
p_3''(x) &= 1.5x > 0 \quad \text{for } x > 0
\end{aligned}
\]
This confirms it's a minimum.

\subsection*{Final Answers}
\begin{enumerate}
\item \( p_3(f; x) = 0.25x^3 - 2.25x + 5 \)
\item \( x_{\min} \approx 1.732 \)
\end{enumerate}

\section*{Problem V: Divided Differences with Repeated Nodes}

\subsection*{Problem Statement}
Consider \( f(x) = x^7 \).

\begin{enumerate}
\item Compute \( f[0, 1, 1, 1, 2, 2] \).
\item We know that this divided difference is expressible in terms of the 5th derivative of \( f \) evaluated at some \( \xi \in (0, 2) \). Determine \( \xi \).
\end{enumerate}

\subsection*{Solution}

\subsubsection*{Part 1: Computing \( f[0, 1, 1, 1, 2, 2] \) by Divided Difference Table}

For \( f(x) = x^7 \), we construct the divided difference table for nodes \( 0, 1, 1, 1, 2, 2 \).

For repeated nodes, we use derivatives:
\[
\begin{aligned}
f[1,1] &= f'(1) = 7 \cdot 1^6 = 7, \\
f[1,1,1] &= \frac{f''(1)}{2!} = \frac{42 \cdot 1^5}{2} = 21, \\
f[2,2] &= f'(2) = 7 \cdot 2^6 = 448
\end{aligned}
\]

Now construct the complete table:

\[
\renewcommand{\arraystretch}{1.5}
\begin{array}{ccccccc}
x & f[] & f[,] & f[,,] & f[,,,] & f[,,,,] & f[,,,,,] \\
\hline
0 & 0 \\
 & & f[0,1] = 1 \\
1 & 1 & & f[0,1,1] = 6 \\
 & & f[1,1] = 7 & & f[0,1,1,1] = 15 \\
1 & 1 & & f[1,1,1] = 21 & & f[0,1,1,1,2] = 42 \\
 & & f[1,1] = 7 & & f[1,1,1,2] = 99 & & f[0,1,1,1,2,2] = 30 \\
1 & 1 & & f[1,1,2] = 120 & & f[1,1,1,2,2] = 102 \\
 & & f[1,2] = 127 & & f[1,1,2,2] = 201 \\
2 & 128 & & f[1,2,2] = 321 \\
 & & f[2,2] = 448 \\
2 & 128 \\
\end{array}
\]

Therefore: \( f[0, 1, 1, 1, 2, 2] = 30 \)

\subsubsection*{Part 2: Determining \( \xi \)}

From the theory of divided differences with repeated nodes:
\[
\begin{aligned}
f[0,1,1,1,2,2] &= \frac{f^{(5)}(\xi)}{5!} \quad \text{for some } \xi \in (0, 2)
\end{aligned}
\]

For \( f(x) = x^7 \), we have:
\[
\begin{aligned}
f^{(5)}(x) &= 7 \cdot 6 \cdot 5 \cdot 4 \cdot 3 \cdot x^2 = 2520x^2
\end{aligned}
\]

So:
\[
\begin{aligned}
30 &= \frac{2520\xi^2}{120}
\end{aligned}
\]

\[
\begin{aligned}
\xi &= \sqrt{\frac{10}{7}} \approx 1.195
\end{aligned}
\]

Since \( \xi \in (0, 2) \), this is valid.

\subsection*{Final Answers}
\begin{enumerate}
\item \( f[0, 1, 1, 1, 2, 2] = 30 \)
\item \( \xi = \sqrt{\frac{10}{7}} \approx 1.195 \)
\end{enumerate}

\section*{Problem VI: Hermite Interpolation and Error Estimation}

\subsection*{Problem Statement}
\( f \) is a function on \([0, 3]\) for which one knows that

\[
f(0) = 1, \, f(1) = 2, \, f'(1) = -1, \, f(3) = f'(3) = 0.
\]

\begin{enumerate}
\item Estimate \( f(2) \) using Hermite interpolation.
\item Estimate the maximum possible error of the above answer if one knows, in addition, that \( f \in C^5[0, 3] \) and \( |f^{(5)}(x)| \leq M \) on \([0, 3]\). Express the answer in terms of \( M \).
\end{enumerate}

\subsection*{Solution}

\subsubsection*{Part 1: Hermite Interpolation to Estimate \( f(2) \)}

We construct the divided difference table for nodes \( 0, 1, 1, 3, 3 \):

\[
\renewcommand{\arraystretch}{1.5}
\begin{array}{cccccc}
x & f[] & \text{1st order} & \text{2nd order} & \text{3rd order} & \text{4th order} \\
\hline
0 & 1 \\
 & & 1 \\
1 & 2 & & -2 \\
 & & -1 & & \frac{2}{3} \\
1 & 2 & & 0 & & -\frac{5}{36}\\
 & & -1 & & \frac{1}{4} \\
3 & 0 & & \frac{1}{2} \\
 & & 0 \\
3 & 0 \\
\end{array}
\]

The Hermite interpolating polynomial is:
\[
\begin{aligned}
H_4(x) = &1 + 1(x-0) - 2(x-0)(x-1) + \frac{2}{3}(x-0)(x-1)^2 - \frac{5}{36}(x-0)(x-1)^2(x-3)
\end{aligned}
\]

Now evaluate at \( x = 2 \):
\[
\begin{aligned}
H_4(2) &= 1 + 1(2) - 2(2)(1) + \frac{2}{3}(2)(1)^2 - \frac{5}{36}(2)(1)^2(-1) \\
&= 1 + 2 - 4 + \frac{4}{3} + \frac{10}{36} \\
&= -1 + \frac{4}{3} + \frac{5}{18} \\
&= -1 + \frac{24}{18} + \frac{5}{18} \\
&= -1 + \frac{29}{18} = \frac{11}{18}
\end{aligned}
\]

Therefore, \( f(2) \approx \frac{11}{18} \).

\subsubsection*{Part 2: Error Estimation}

The error formula for Hermite interpolation is:
\[
\begin{aligned}
f(x) - H_4(x) &= \frac{f^{(5)}(\xi)}{5!} \prod_{i=0}^{2} (x - x_i)^{m_i} \\
&= \frac{f^{(5)}(\xi)}{120} (x-0)^1(x-1)^2(x-3)^2
\end{aligned}
\]

For \( x = 2 \):
\[
\begin{aligned}
|f(2) - H_4(2)| &= \left| \frac{f^{(5)}(\xi)}{120} (2-0)(2-1)^2(2-3)^2 \right| \\
&= \left| \frac{f^{(5)}(\xi)}{120} \cdot 2 \cdot 1 \cdot 1 \right| \\
&= \left| \frac{f^{(5)}(\xi)}{60} \right| \\
&\leq \frac{M}{60}
\end{aligned}
\]

\subsection*{Final Answers}
\begin{enumerate}
\item \( f(2) \approx \frac{11}{18} \).
\item Maximum possible error: \( \frac{M}{60} \)
\end{enumerate}


\section*{Problem VII: Relationship Between Differences and Divided Differences}

\subsection*{Problem Statement}
Define forward difference by

\[
\Delta f(x) = f(x + h) - f(x),
\]
\[
\Delta^{k+1}f(x) = \Delta\Delta^k f(x) = \Delta^k f(x + h) - \Delta^k f(x)
\]

and backward difference by

\[
\nabla f(x) = f(x) - f(x - h),
\]
\[
\nabla^{k+1}f(x) = \nabla\nabla^k f(x) = \nabla^k f(x) - \nabla^k f(x - h).
\]

For \( x_j = x + jh \), prove

\[
\Delta^k f(x) = k!h^k f[x_0, x_1, \ldots, x_k],
\]
\[
\nabla^k f(x) = k!h^k f[x_0, x_{-1}, \ldots, x_{-k}].
\]

\subsection*{Proof}

\subsubsection*{Part 1: Forward Difference}

Let \( x_j = x + jh \) for \( j = 0, 1, \ldots, k \). We prove by induction that:

\[
\Delta^k f(x) = k!h^k f[x_0, x_1, \ldots, x_k]
\]

\textbf{Base case: } \( k = 1 \)

\[
\begin{aligned}
\Delta f(x) &= f(x + h) - f(x) \\
&= \frac{f(x + h) - f(x)}{h} \cdot h \\
&= f[x_0, x_1] \cdot h \\
&= 1!h^1 f[x_0, x_1]
\end{aligned}
\]

So the formula holds for \( k = 1 \).

\textbf{Inductive step: }

Assume the formula holds for \( k \):

\[
\Delta^k f(x) = k!h^k f[x_0, x_1, \ldots, x_k]
\]

Now for \( k+1 \):

\[
\begin{aligned}
\Delta^{k+1} f(x) &= \Delta^k f(x + h) - \Delta^k f(x) \\
&= k!h^k f[x_1, x_2, \ldots, x_{k+1}] - k!h^k f[x_0, x_1, \ldots, x_k] \\
&= k!h^k \left( f[x_1, \ldots, x_{k+1}] - f[x_0, \ldots, x_k] \right)
\end{aligned}
\]

Using the recursive property of divided differences:

\[
\begin{aligned}
f[x_0, x_1, \ldots, x_{k+1}] &= \frac{f[x_1, \ldots, x_{k+1}] - f[x_0, \ldots, x_k]}{x_{k+1} - x_0} \\
&= \frac{f[x_1, \ldots, x_{k+1}] - f[x_0, \ldots, x_k]}{(k+1)h}
\end{aligned}
\]

Therefore:

\[
\begin{aligned}
\Delta^{k+1} f(x) &= k!h^k \cdot (k+1)h \cdot f[x_0, x_1, \ldots, x_{k+1}] \\
&= (k+1)!h^{k+1} f[x_0, x_1, \ldots, x_{k+1}]
\end{aligned}
\]

This completes the inductive step.

\subsubsection*{Part 2: Backward Difference}

Let \( x_j = x + jh \) for \( j = 0, -1, \ldots, -k \). We prove by induction that:

\[
\nabla^k f(x) = k!h^k f[x_0, x_{-1}, \ldots, x_{-k}]
\]

\textbf{Base case: } \( k = 1 \)

\[
\begin{aligned}
\nabla f(x) &= f(x) - f(x - h) \\
&= \frac{f(x) - f(x - h)}{h} \cdot h \\
&= f[x_0, x_{-1}] \cdot h \\
&= 1!h^1 f[x_0, x_{-1}]
\end{aligned}
\]

So the formula holds for \( k = 1 \).

\textbf{Inductive step: }

Assume the formula holds for \( k \):

\[
\nabla^k f(x) = k!h^k f[x_0, x_{-1}, \ldots, x_{-k}]
\]

Now for \( k+1 \):

\[
\begin{aligned}
\nabla^{k+1} f(x) &= \nabla^k f(x) - \nabla^k f(x - h) \\
&= k!h^k \left( f[x_0, x_{-1}, \ldots, x_{-k}] - f[x_{-1}, x_{-2}, \ldots, x_{-(k+1)}] \right)
\end{aligned}
\]

Using the recursive property of divided differences:

\[
\begin{aligned}
f[x_{-1}, x_0, x_{-1}, \ldots, x_{-k}] &= \frac{f[x_0, x_{-1}, \ldots, x_{-k}] - f[x_{-1}, x_{-2}, \ldots, x_{-(k+1)}]}{x_0 - x_{-(k+1)}} \\
&= \frac{f[x_0, x_{-1}, \ldots, x_{-k}] - f[x_{-1}, x_{-2}, \ldots, x_{-(k+1)}]}{(k+1)h}
\end{aligned}
\]

Therefore:

\[
\begin{aligned}
\nabla^{k+1} f(x) &= k!h^k \cdot (k+1)h \cdot f[x_0, x_{-1}, \ldots, x_{-(k+1)}] \\
&= (k+1)!h^{k+1} f[x_0, x_{-1}, \ldots, x_{-(k+1)}]
\end{aligned}
\]

This completes the inductive step.

\subsection*{Conclusion}
We have proved both identities using mathematical induction:
\[
\begin{aligned}
\Delta^k f(x) &= k!h^k f[x_0, x_1, \ldots, x_k]
\end{aligned}
\]

\[
\begin{aligned}
\nabla^k f(x) &= k!h^k f[x_0, x_{-1}, \ldots, x_{-k}]
\end{aligned}
\]


\section*{Problem VIII: Partial Derivatives of Divided Differences}

\subsection*{Problem Statement}
Assume \( f \) is differentiable at \( x_0 \). Prove

\[
\frac{\partial}{\partial x_0} f[x_0, x_1, \ldots, x_n] = f[x_0, x_0, x_1, \ldots, x_n].
\]

What about the partial derivative with respect to one of the other variables?

\subsection*{Proof}

\subsubsection*{Part 1: Proof for \( \frac{\partial}{\partial x_0} \)}

We prove the identity by induction on \( n \).

\textbf{Base case: } \( n = 0 \)

For \( n = 0 \), the divided difference is just the function value:
\[
f[x_0] = f(x_0)
\]

Then:
\[
\frac{\partial}{\partial x_0} f[x_0] = f'(x_0)
\]

By definition of divided differences with repeated nodes:
\[
f[x_0, x_0] = f'(x_0)
\]

So the identity holds for \( n = 0 \).

\textbf{Inductive step: }

Assume the identity holds for \( n-1 \):
\[
\frac{\partial}{\partial x_0} f[x_0, x_1, \ldots, x_{n-1}] = f[x_0, x_0, x_1, \ldots, x_{n-1}]
\]

Now consider \( f[x_0, x_1, \ldots, x_n] \). Using the recursive definition of divided differences:
\[
f[x_0, x_1, \ldots, x_n] = \frac{f[x_1, \ldots, x_n] - f[x_0, \ldots, x_{n-1}]}{x_n - x_0}
\]

Taking partial derivative with respect to \( x_0 \):
\[
\begin{aligned}
\frac{\partial}{\partial x_0} f[x_0, x_1, \ldots, x_n] 
&= \frac{\partial}{\partial x_0} \left( \frac{f[x_1, \ldots, x_n] - f[x_0, \ldots, x_{n-1}]}{x_n - x_0} \right)
\end{aligned}
\]

\[
\begin{aligned}
&= \frac{ - \frac{\partial}{\partial x_0} f[x_0, \ldots, x_{n-1}] \cdot (x_n - x_0) + (f[x_1, \ldots, x_n] - f[x_0, \ldots, x_{n-1}]) }{(x_n - x_0)^2}
\end{aligned}
\]

Using the inductive hypothesis:
\[
\frac{\partial}{\partial x_0} f[x_0, \ldots, x_{n-1}] = f[x_0, x_0, x_1, \ldots, x_{n-1}]
\]

So:
\[
\begin{aligned}
\frac{\partial}{\partial x_0} f[x_0, x_1, \ldots, x_n] 
&= \frac{ - f[x_0, x_0, x_1, \ldots, x_{n-1}] \cdot (x_n - x_0) + (f[x_1, \ldots, x_n] - f[x_0, \ldots, x_{n-1}]) }{(x_n - x_0)^2}
\end{aligned}
\]

Now, by the recursive definition of divided differences with repeated nodes:
\[
f[x_0, x_0, x_1, \ldots, x_n] = \frac{f[x_0, x_1, \ldots, x_n] - f[x_0, x_0, \ldots, x_{n-1}]}{x_n - x_0}
\]

Rearranging this:
\[
f[x_0, x_0, x_1, \ldots, x_n] = \frac{f[x_0, x_1, \ldots, x_n] - f[x_0, x_0, \ldots, x_{n-1}]}{x_n - x_0}
\]

But from our expression above, we have exactly:
\[
\frac{\partial}{\partial x_0} f[x_0, x_1, \ldots, x_n] = f[x_0, x_0, x_1, \ldots, x_n]
\]

This completes the inductive step.

\subsubsection*{Part 2: Partial Derivatives with Respect to Other Variables}

The key insight is that divided differences are \textbf{symmetric functions} of their nodes:
\[
f[x_0, x_1, \ldots, x_n] = f[x_{\sigma(0)}, x_{\sigma(1)}, \ldots, x_{\sigma(n)}]
\]
for any permutation $\sigma$ of $\{0, 1, \ldots, n\}$.

Now, for any $0 \leq k \leq n$, we can prove:
\[
\frac{\partial}{\partial x_k} f[x_0, x_1, \ldots, x_n] = f[x_0, \ldots, x_{k-1}, x_k, x_k, x_{k+1}, \ldots, x_n]
\]

\textbf{Proof:}

Choose a permutation $\sigma$ such that $\sigma(0) = k$. Then:
\[
\begin{aligned}
\frac{\partial}{\partial x_k} f[x_0, x_1, \ldots, x_n] 
&= \frac{\partial}{\partial x_k} f[x_{\sigma(0)}, x_{\sigma(1)}, \ldots, x_{\sigma(n)}] \quad \text{(by symmetry)} \\

&= f[x_{\sigma(0)}, x_{\sigma(0)}, x_{\sigma(1)}, \ldots, x_{\sigma(n)}]  \\

&= f[x_k, x_k, x_{\sigma(1)}, \ldots, x_{\sigma(n)}]
\end{aligned}
\]

Since $\sigma(1), \ldots, \sigma(n)$ is just a reordering of $\{x_0, \ldots, x_n\} \setminus \{x_k\}$, we can rewrite this as:
\[
f[x_0, \ldots, x_{k-1}, x_k, x_k, x_{k+1}, \ldots, x_n]
\]

This completes the proof for all $k = 0, 1, \ldots, n$.


\subsection*{Problem IX. A Min-Max Problem}

For \( n \in \mathbb{N}^+ \) and a fixed \( a_0 \neq 0 \), determine
\[
\min \max_{x \in [a,b]} |a_0 x^n + a_1 x^{n-1} + \cdots + a_n|
\]
where the minimum is taken over all \( a_i \in \mathbb{R}, \, i = 1, 2, \cdots, n \).

\subsubsection*{Solution}

\paragraph{Case Analysis:}

\textbf{Case 1: Interval \( [-1,1] \)}

For the special case \( [a,b] = [-1,1] \), the solution is given by the Chebyshev polynomial of degree \( n \):
\[
T_n(x) = \cos(n \arccos x)
\]
The scaled Chebyshev polynomial:
\[
P^*(x) = \frac{a_0}{2^{n-1}} T_n(x)
\]
achieves the minimum maximum deviation:
\[
\min \max_{x \in [-1,1]} |P(x)| = \frac{|a_0|}{2^{n-1}}
\]

\textbf{Case 2: General interval \( [a,b] \)}

For a general interval \( [a,b] \), we use the linear transformation:
\[
x = \frac{b-a}{2}t + \frac{a+b}{2}
\]
which maps \( t \in [-1,1] \) to \( x \in [a,b] \).

The optimal polynomial is:
\[
P^*(x) = \frac{a_0}{(b-a)^n} \cdot 2^{2n-1} \cdot T_n\left( \frac{2x - (a+b)}{b-a} \right)
\]
and the minimal maximum value is:
\[
\min \max_{x \in [a,b]} |P(x)| = \frac{|a_0|(b-a)^n}{2^{2n-1}}
\]


\subsection*{Problem X. Imitate the proof of Chebyshev Theorem}

Express the Chebyshev polynomial of degree \( n \in \mathbb{N} \) as a polynomial \( T_n \) and consider its domain extended from \([-1, 1]\) to \(\mathbb{R}\). For a fixed \( a > 1 \), define
\[
\mathbb{P}_n^a := \{ p \in \mathbb{P}_n : p(a) = 1 \}
\]
and a specific polynomial within this set:
\[
\hat{p}_n(x) := \frac{T_n(x)}{T_n(a)}.
\]
For the uniform norm defined by
\[
\|f\|_\infty = \max_{x \in [-1, 1]} |f(x)|,
\]
prove that \(\hat{p}_n\) is the minimizer of this norm over the set \(\mathbb{P}_n^a\), i.e.,
\[
\forall p \in \mathbb{P}_n^a, \quad \|\hat{p}_n\|_\infty \leq \|p\|_\infty.
\]

\subsubsection*{Proof}

Let \( \hat{p}_n(x) = \frac{T_n(x)}{T_n(a)} \). Note that since \( a > 1 \) and the leading coefficient of \( T_n \) is positive for \( n \geq 0 \), we have \( T_n(a) > 0 \). Thus, \( \hat{p}_n \in \mathbb{P}_n^a \) because \( \hat{p}_n(a) = \frac{T_n(a)}{T_n(a)} = 1 \).

Recall the key properties of the Chebyshev polynomial \( T_n(x) = \cos(n \arccos x) \) on the interval \([-1, 1]$:
\begin{itemize}
    \item It satisfies \( |T_n(x)| \leq 1 \) for all \( x \in [-1, 1] \).
    \item It attains its extreme values \( \pm 1 \) at \( n+1 \) distinct points \( x_k = \cos\left(\frac{k\pi}{n}\right) \), for \( k = 0, 1, \dots, n \), in the interval \([-1, 1]$, with alternating signs. That is,
    \[
    T_n(x_k) = (-1)^k.
    \]
\end{itemize}

Consequently, for \( \hat{p}_n \), we have:
\[
|\hat{p}_n(x)| = \left| \frac{T_n(x)}{T_n(a)} \right| \leq \frac{1}{T_n(a)} \quad \text{for all } x \in [-1, 1],
\]
and at the points \( x_k \),
\[
\hat{p}_n(x_k) = \frac{(-1)^k}{T_n(a)}.
\]
Thus, \( \hat{p}_n \) oscillates between \( \frac{1}{T_n(a)} \) and \( -\frac{1}{T_n(a)} \) at the \( n+1 \) points \( x_k \), and \( \|\hat{p}_n\|_\infty = \frac{1}{T_n(a)} \).

Now, suppose for the sake of contradiction that there exists a polynomial \( q \in \mathbb{P}_n^a \) such that
\[
\|q\|_\infty < \|\hat{p}_n\|_\infty = \frac{1}{T_n(a)}.
\]
Define the difference polynomial:
\[
r(x) = \hat{p}_n(x) - q(x).
\]
Note that \( r(x) \) is a polynomial of degree at most \( n \). We will examine the sign of \( r(x) \) at the alternation points \( x_k \) of \( \hat{p}_n \).

Since \( \hat{p}_n, q \in \mathbb{P}_n^a \), we have \( \hat{p}_n(a) = q(a) = 1 \), which implies:
\[
r(a) = \hat{p}_n(a) - q(a) = 1 - 1 = 0.
\]
This is crucial and will be used later.

By our assumption \( \|q\|_\infty < \frac{1}{T_n(a)} \), we have \( |q(x)| < \frac{1}{T_n(a)} \) for all \( x \in [-1, 1] \).

Now consider the values of \( r(x) \) at the points \( x_k \):
\begin{align*}
\text{For even } k: \quad & \hat{p}_n(x_k) = \frac{1}{T_n(a)}, \quad \text{so } r(x_k) = \frac{1}{T_n(a)} - q(x_k) > 0. \\
\text{For odd } k: \quad & \hat{p}_n(x_k) = -\frac{1}{T_n(a)}, \quad \text{so } r(x_k) = -\frac{1}{T_n(a)} - q(x_k) < 0.
\end{align*}
Therefore, the polynomial \( r(x) \) alternates in sign at the \( n+1 \) points \( x_0, x_1, \dots, x_n \). By the Intermediate Value Theorem, \( r(x) \) must have at least \( n \) distinct roots in the open interval \( (-1, 1) \), one between each consecutive pair \( (x_k, x_{k+1}) \).

Additionally, we have already established that \( r(a) = 0 \), and since \( a > 1 \), this root lies outside \( [-1, 1] \) and is distinct from the roots found inside \( (-1, 1) \).

Thus, in total, the polynomial \( r(x) \), which has degree at most \( n \), is found to have at least \( n \) roots (in \((-1,1)\)) + 1 root (at \(x=a\)) = \( n+1 \) distinct roots. The only polynomial of degree at most \( n \) with \( n+1 \) or more distinct roots is the zero polynomial. Therefore, \( r(x) \equiv 0 \), implying \( q(x) \equiv \hat{p}_n(x) \).

But this contradicts our assumption that \( \|q\|_\infty < \|\hat{p}_n\|_\infty \). Hence, our initial assumption must be false.

We conclude that no such polynomial \( q \) exists, and therefore,
\[
\forall p \in \mathbb{P}_n^a, \quad \|\hat{p}_n\|_\infty \leq \|p\|_\infty.
\]
This completes the proof.


\subsection*{Proof of Lemma 2.53}

We want to prove that for Bernstein basis polynomials:
\[
b_{n-1,k}(t) = \frac{n - k}{n} b_{n,k}(t) + \frac{k + 1}{n} b_{n,k+1}(t)
\]
where the Bernstein basis polynomial of degree $n$ is defined as:
\[
b_{n,k}(t) = \binom{n}{k} t^k (1-t)^{n-k}, \quad \text{for } k = 0, 1, \ldots, n.
\]

\subsubsection*{Proof}

We start from the right-hand side and show that it equals the left-hand side:
\[
\begin{aligned}
&\frac{n - k}{n} b_{n,k}(t) + \frac{k + 1}{n} b_{n,k+1}(t) \\
&= \frac{n - k}{n} \binom{n}{k} t^k (1-t)^{n-k} + \frac{k + 1}{n} \binom{n}{k+1} t^{k+1} (1-t)^{n-k-1}.
\end{aligned}
\]

Using the identities for binomial coefficients:
\[
\binom{n}{k} = \frac{n!}{k!(n-k)!} \quad \text{and} \quad \binom{n}{k+1} = \frac{n!}{(k+1)!(n-k-1)!},
\]
we substitute and simplify:
\[
\begin{aligned}
&\frac{n - k}{n} \cdot \frac{n!}{k!(n-k)!} t^k (1-t)^{n-k} + \frac{k + 1}{n} \cdot \frac{n!}{(k+1)!(n-k-1)!} t^{k+1} (1-t)^{n-k-1} \\
&= \frac{n!}{n \cdot k!(n-k-1)!} t^k (1-t)^{n-k} + \frac{n!}{n \cdot k!(n-k-1)!} t^{k+1} (1-t)^{n-k-1} \\
&= \frac{n!}{n \cdot k!(n-k-1)!} \left[ t^k (1-t)^{n-k} + t^{k+1} (1-t)^{n-k-1} \right].
\end{aligned}
\]

Factoring out common terms:
\[
\begin{aligned}
&= \frac{n!}{n \cdot k!(n-k-1)!} \cdot t^k (1-t)^{n-k-1} \left[ (1-t) + t \right] \\
&= \frac{n!}{n \cdot k!(n-k-1)!} \cdot t^k (1-t)^{n-k-1}.
\end{aligned}
\]

Now, recall the definition:
\[
b_{n-1,k}(t) = \binom{n-1}{k} t^k (1-t)^{n-1-k} = \frac{(n-1)!}{k!(n-k-1)!} t^k (1-t)^{n-k-1}.
\]

Since:
\[
\frac{n!}{n \cdot k!(n-k-1)!} = \frac{(n-1)!}{k!(n-k-1)!},
\]
we conclude that:
\[
\frac{n - k}{n} b_{n,k}(t) + \frac{k + 1}{n} b_{n,k+1}(t) = b_{n-1,k}(t),
\]
which completes the proof.


\subsection*{Proof of Lemma 2.55}

We want to prove that for Bernstein basis polynomials:
\[
\int_{0}^{1} b_{n,k}(t) \, dt = \frac{1}{n+1}, \quad \forall k = 0, 1, \ldots, n
\]
where the Bernstein basis polynomial of degree $n$ is defined as:
\[
b_{n,k}(t) = \binom{n}{k} t^k (1-t)^{n-k}.
\]

\subsubsection*{Proof}

We compute the integral directly using the Beta function. Recall that the Beta function is defined as:
\[
B(x,y) = \int_0^1 t^{x-1} (1-t)^{y-1} \, dt = \frac{\Gamma(x)\Gamma(y)}{\Gamma(x+y)}, \quad \text{for } x > 0, y > 0.
\]

For our integral:
\begin{align*}
\int_{0}^{1} b_{n,k}(t) \, dt &= \int_{0}^{1} \binom{n}{k} t^k (1-t)^{n-k} \, dt \\
&= \binom{n}{k} \int_{0}^{1} t^{k} (1-t)^{n-k} \, dt \\
&= \binom{n}{k} B(k+1, n-k+1).
\end{align*}

Using the relationship between the Beta and Gamma functions:
\[
B(k+1, n-k+1) = \frac{\Gamma(k+1)\Gamma(n-k+1)}{\Gamma(n+2)}.
\]

Since $k$ and $n-k$ are non-negative integers, we can use the factorial property $\Gamma(m+1) = m!$:
\[
B(k+1, n-k+1) = \frac{k! \cdot (n-k)!}{(n+1)!}.
\]

Now substitute back:
\begin{align*}
\int_{0}^{1} b_{n,k}(t) \, dt &= \binom{n}{k} \cdot \frac{k! \cdot (n-k)!}{(n+1)!} \\
&= \frac{n!}{k!(n-k)!} \cdot \frac{k! \cdot (n-k)!}{(n+1)!} \\
&= \frac{n!}{(n+1)!} \\
&= \frac{1}{n+1}.
\end{align*}


\end{document}